\documentstyle[% 


righttag% 


,11pt% 


,amssymb% 


,russian%         remove in Eng. 


,izvru%           Set izven% in Eng. 


% ,izvru-ks       for brief communications 


]{amsart} 


 


%    Math definitions 


 


\newcommand{\col}{\operatorname{col}} 


\newcommand{\vraisup}{\operatornamewithlimits{vrai\,sup}} 


\newcommand{\vraiinf}{\operatornamewithlimits{vrai\,inf}} 


 


\theoremstyle{plain} 


\theorembodyfont{\it} 


\newtheorem{thmnonum}{\hskip\parindent Теорема} 


\renewcommand{\thethmnonum}{} 


\newtheorem{lemnonum}{\hskip\parindent Лемма} 


\renewcommand{\thelemnonum}{} 


 


\theoremstyle{definition} 


\theorembodyfont{\rm} 


\newtheorem{notenonum}{\hskip\parindent Замечание} 


\renewcommand{\thenotenonum}{} 





\renewcommand{\baselinestretch}{1.1}


 


%\hoffset=-15mm            For Eng. 


 


\setcounter{page}{37} 


\god{1997} 


\nomer{6~(421)} %          Remove in Eng. 


%\nomer{Vol.\,40, No.\,8}  Set in Eng. 


% \str{75-81}              Set in Eng. 


\udk{517.929} 


\author{\it П.М.\,СИМОНОВ, А.В.\,ЧИСТЯКОВ} 


\title{Об экспоненциальной устойчивости линейных\\


дифференциально-разностных систем}





\thanks{Работа выполнена при финансовой поддержке Российского фонда


фундаментальных исследований (грант 96-01-01613).}


\date{} 


% \pagestyle{myheadings}   Set in Eng. 


 


\pagestyle{plain} %    Renove in Eng. 


\begin{document} 


 


% \thispagestyle{plain}   Set in Eng. 


 


\maketitle 


 


% Body text 


 





\section{Теорема об экспоненциальной устойчивости}





Признаки экспоненциальной устойчивости линейных


дифференциально-разностных уравнений (ЛДРУ) в явном виде содержат


условие устойчивости ``нейтральной части'' --- обратимость разностного


оператора, действующего на старшую производную (\cite{Hale},


п.\,12.4--12.6; \cite{Kol},


гл.\,3).


Вопрос о необходимости этого условия изучал В.Г.Курбатов в работах


\cite{Ku1}, \cite{Ku2}. В (\cite{Ku3}, гл.\,5)


показано, что устойчивость нейтральной части является


необходимым условием экспоненциального убывания решения и его


производной по норме $\bold L_p(1\le p\le\infty)$. Такое свойство


решений ЛДРУ называется в \cite{Ku3} экспоненциальной устойчивостью ЛДРУ по


норме $\bold W_p^{(1)}$. Отметим, что согласно теореме 6.3.11


(\cite{Ku3}, сс.\,111,~112)


для ЛДРУ с гладкими коэффициентами $\bold W_p^{(1)}$-устойчивость


эквивалентна устойчивости по Ляпунову (т.е. устойчивости в чебышевской


метрике).





В предлагаемой работе сформулирована теорема, дополняющая результаты


В.Г.\,Курбатова. В этой теореме утверждается, в частности, что при


начальных функциях класса $\bold W_1^{(1)}$ экспоненциальная


устойчивость по Ляпунову влечет устойчивость нейтральной части в


$\bold L_1$-метрике. Отсюда и из результатов \cite{Ku3} следует, что


экспоненциальная устойчивость эквивалентна 


$\bold W_1^{(1)}$-устойчивости при постоянно действующих 


$\bold L_1$-возмущениях. Результаты работы были анонсированы ранее


в \cite{Si1}, \cite{Si2}.





В статье используются следующие определения и обозначения.


$\bold R^n$ -- евклидово пространство $n$-мерных вектор-столбцов


$x=\col\{x_1,\dotsc,x_n\}$ со скалярным произведением


$\langle x,y\rangle=\sum\limits_{i=1}^nx_iy_i$,


нормой $\|x\|=\sqrt{\smash[b]{\langle x,x\rangle}}$


и координатным базисом $\{e_i\}_{i=1}^n$, причем 


$E=\{e_1,\dotsc,e_n\}$ --- единичная $n\times n$-матрица.


$\bold M^n$ ---


пространство $n\times n$-матриц $A=\{a_{ij}\}_{i,j=1}^n$ с операторной


нормой $\|A\|$. $\bold L_1(\Delta)$ --- пространство (классов


эквивалентности) измеримых функций $z:\Delta\to\bold R^n$,


суммируемых по мере Лебега на промежутке $\Delta\subset\bold R$, с


нормой $\|z\|_{\bold L_1(\Delta)}=\int\limits_\Delta\|z(t)\|dt$;


$\bold L_\infty(\Delta)$ --- пространство (классов


эквивалентности) вектор-функций $y:\Delta\to\bold R^n$, измеримых


и ограниченных в существенном на промежутке $\Delta\subset\bold R$,


с нормой


$\|z\|_{\bold L_\infty(\Delta)}=\vraisup\limits_{t\in\Delta}\|z(t)\|$;


$\bold W_1(\Delta)=\bold W_1^{(1)}(\Delta)$ --- пространство


функций $x:\Delta\to\bold R^n$, абсолютно непрерывных на


промежутке $\Delta\subset\bold R$ и ограниченных по норме


$\|x\|_{\bold W_1(\Delta)}=\|x\|_{\bold L_1(\Delta)}+\|\Dot


x\|_{\bold L_1(\Delta)}$.





Через $\bold B^*$ обозначим пространство, сопряженное линейному


нормированному пространству $\bold B$, а через 


$Q^*:\bold B^*\to\bold B^*$ --- оператор, сопряженный линейному


ограниченному оператору $Q:\bold B\to\bold B$.


Пространство $(\bold L_1(\Delta))^*$


отождествляем с пространством $\bold L_\infty(\Delta)$


посредством канонической двойственности


%


$$


(z,y)=\int_\Delta\langle z(t),y(t)\rangle dt,


\quad z\in\bold L_1(\Delta),\quad y\in\bold


L_\infty(\Delta).


$$





Всюду ниже $\chi_\Delta$ --- характеристическая функция множества


$\Delta$.





Рассмотрим уравнение


%


\begin{equation}


\Dot x(t)=\sum_{g\in G_1}A_g(t)\Dot x(t-g)+\sum_{h\in


H}B_h(t)x(t-h),\quad t\in\bold R,


\end{equation}


%


где $G_1$ --- конечное множество чисел $g>0$, $H$ --- конечное множество


чисел $h\ge0$, матрицы-функции $A_g:\bold R\to\bold M^n$


и


$B_h:\bold R\to\bold M^n$ измеримы по Лебегу и удовлетворяют условиям


% ---------------------------------------------  v


$$


\vraisup_{t\in\bold R}\|A_g(t)\|<\infty\quad\text{и}\quad


\sup_{t\in\bold R}


\int_t^{t+1}\|B_h(s)\|ds<\infty


$$


при всех $g\in G_1$ и $h\in H$.





В работах \cite{Hale}--\cite{Ku3}


ставится для уравнения (1) следующая начальная задача.


Пусть заданы начальный момент $t_0\ge 0$ и начальная функция 


$\phi:(-\infty,t_0]\to\bold R^n$. Предполагается, что функция $\phi$ абсолютно


непрерывна на каждом отрезке полуоси $(-\infty,t_0]$. Требуется найти


функцию $x:\bold R\to\bold R^n$,


удовлетворяющую равенству (1) при


почти всех $t\ge t_0$ и


начальному условию


% ------------------- v


\begin{equation}


x(t)=\phi(t)\text{ при всех }t\le t_0.


\end{equation}


%


Хорошо известно (\cite{Kol}, с.\,42; \cite{Ku3}, с.\,82),


что при любых начальных


данных $t_0$ и $\phi$ задача (1), (2) имеет единственное


решение $x(t)=x(t,t_0,\phi)$.





Различные определения устойчивости связаны с выбором разных метрик в


пространстве начальных функций и пространстве решений (\cite{Hale},


с.\,130; \cite{Kol},


с.\,105; \cite{Ku3}, с.\,90). В данной статье принято





\begin{defn} Будем говорить, что уравнение (1) равномерно (по


начальному моменту $t_0\ge0$) экспоненциально устойчиво, если


существуют такие константы $N>0$ и $\gamma>0$, что при всех $t_0\ge0$


и $\phi\in\bold W_1(-\infty,t_0]$ решение $x=x(\cdot,t_0,\phi)$


начальной задачи (1), (2) при всех $t\ge t_0$ удовлетворяет


экспоненциальной оценке


%


$$


\|x(t)\|\le Ne^{-\gamma(t-t_0)}\|\phi\|_{\bold W_1(-\infty,t_0]}.


$$


\end{defn}





Иными словами, уравнение (1) равномерно экспоненциально устойчиво,


если тривиальное решение $x(t)=x(t,t_0,0)\equiv 0$ этого уравнения


равномерно экспоненциально устойчиво по Ляпунову относительно


$\bold W_1$-возмущений начальной функции $\phi$.





Наше исследование экспоненциальной устойчивости существенно основано


на специальном представлении решений задачи (1), (2). Как следует из


результатов (\cite{Kol}, гл.\,1; \cite{Ku1}, гл.\,5;


\cite{Az1}, гл.\,5; \cite{Az2}),


это решение представимо в виде формулы Коши


%


\begin{equation}


x(t)=X(t,t_0)\phi(t_0)+\int_{t_0}^tC(t,s)f_\phi(s)ds,\quad t\ge


t_0\ge0.


\end{equation}


%


Здесь $X(t,t_0)$ --- фундаментальная матрица, $C(t,s)$ --- матрица Коши,


функция $f_\phi$ определяется равенством


$$


f_\phi(t)=\sum


\begin{Sb}


g\in G_1: \\ g\ge t-t_0


\end{Sb}


A_g(t)\Dot\phi(t-g)+\sum


\begin{Sb}


h\in H: \\ h\ge t-t_0


\end{Sb}


B_h(t)\phi(t-h).


$$


Фундаментальная матрица $X_{t_0}(t)\equiv X(t,t_0)$ является


(\cite{Az1}, гл.\,3, \S\,3.4)


решением матричной задачи Коши


%


\begin{equation}


\begin{cases}


\Dot X_{t_0}(t)=\sum\limits_{g\in G_1}A_g(t)\Dot


X_{t_0}(t-g)+\sum\limits_{h\in


H}B_h(t)X_{t_0}(t-h) &\text{ при }t\ge t_0\ge 0,\\


%


X_{t_0}(t)=0 &\text{ при }t<t_0,\quad X_{t_0}(t_0)=E.


\end{cases}


\end{equation}


%


Матрица Коши $C(t,s)$ есть ядро интегрального представления решения


полуоднородной задачи Коши. Эта задача изучена ниже.





Из формулы Коши (3) следует, что наличие того или иного свойства


устойчивости решений уравнения (1) определяется асимптотическим


поведением фундаментальной матрицы и матрицы Коши. В частности, если


справедливы экспоненциальные оценки


%


\begin{align}


\begin{split}


\|X(t,t_0)\| &\le Ne^{-\gamma(t-t_0)}, \quad t\ge t_0\ge0,\\


%


\|C(t,s)\| &\le Ne^{-\gamma(t-s)}, \quad t\ge s\ge t_0\ge0,


\end{split}


\end{align}


%


где $N,\gamma>0$, то уравнение (1) равномерно экспоненциально


устойчиво. В основной теореме статьи утверждается, что наличие таких


оценок


является не только достаточным, но и необходимым условием


экспоненциальной устойчивости для рассматриваемого класса уравнений.





Для получения условий экспоненциальной устойчивости применяются


различные методы. Наиболее распространенные методы основаны на


построении функций Ляпунова или функционалов Ляпунова--Красовского


\cite{Hale}, \cite{Kol}. Для стационарных уравнений эффективен


операционный метод \cite{Hale}, \cite{Kol}, для периодических --- метод


производящих функций \cite{Ber}--\cite{Si4}.





Мощным методом получения экспоненциальных оценок является $W$-метод


(\cite{Az1}--\cite{Ber}). Суть


метода состоит в редукции задачи об устойчивости к задаче об


обратимости линейного оператора в банаховом пространстве. Такая


редукция основана на обобщениях известной теоремы Боля--Перрона (\cite{Ku3},


гл.\,5; \cite{Az1}, гл.\,5; \cite{Az2}).





Факт экспоненциальной устойчивости уравнения (1) существенно связан с


поведением решений разностного уравнения


%


\begin{equation}


\begin{cases}


y(t)=\sum\limits_{g\in G_1}A_g(t)y(t-g)+p(t), & t\ge0,\\


%


y(\xi)=0, & \text{если }\xi<0,


\end{cases}


\end{equation}


%


при $p\in\bold L_1[0,+\infty)$. Решение $y$ этого уравнения имеет


представление


%


\begin{equation}


y(t)=p(t)+\sum_{g\in G}P_g(t)p(t-g).


\end{equation}


%


Здесь $G$ --- множество всевозможных сумм $g=g_0+\dotsb+g_{m-1}$ отклонений


$g_i\in G_1$, коэффициенты $P_g$ вычисляются по формуле


%


\begin{equation}


P_g(t)=\sum_{g_0+\dotsb+


g_{m-1}=g}\prod_{k=0}^{m-1}A_{g_k}\biggl(t-\sum_{i=0}^kg_i\biggr).


\end{equation}


%





\begin{defn} Будем говорить, что уравнение (6) устойчиво, если


найдется такая положительная константа $q$, что при каждом $b>0$ и для


всех $p\in\bold L_1[0,+\infty)$ решение $y$ уравнения (6)


удовлетворяет оценке


%


\begin{equation}


\int_0^b\|y(t)\|dt\le q\|p\|_{\bold L_1[0,+\infty)}.


\end{equation}


%


\end{defn}





Устойчивость уравнения (6) означает его разрешимость в пространстве


$\bold L_1[0,+\infty)$ или, что эквивалентно,


обратимость разностного оператора 


${\cal I-S}:\bold L_1[0,+\infty)\to\bold L_1[0,+\infty)$,


действующего на производную $\Dot x$


в уравнении (1). Здесь ${\cal I}$ -- тождественный оператор,


\begin{align*}


({\cal S}y)(t) &=\sum_{g\in G_1}A_g(t)({\cal S}_gy)(t),\quad t\ge0,\\


({\cal S}_gy)(t) &=


\begin{cases}


y(t-g), & t\ge g,\\


0, & t<g.


\end{cases}


\end{align*}





Разрешив уравнение (1) относительно $\Dot x$, получим уравнение


%


\begin{equation}


\Dot x(t)=\sum_{g\in G}\sum_{h\in H}R_{gh}(t)x(t-g-h),\quad t\in\bold R,


\end{equation}


%


где $R_{gh}(t)=P_g(t)B_h(t-g)$. Это уравнение можно дополнить


начальным условием типа (2). Однако следует отметить, что при этом


множества решений полученной задачи и исходной задачи (1), (2) не


совпадают. Поэтому асимптотические свойства решений уравнения (10)


существенно отличаются от асимптотических свойств решений уравнения


(1).





Начальная задача (1), (2) сводится к задаче Коши


%


\begin{equation}


\begin{cases}


\Dot x(t)=\sum\limits_{g\in G}\sum\limits_{h\in H}R_{gh}(t)


x(t-g-h)+f(t), &t\ge t_0\ge0,\\


%


x(\xi)=0, & \text{при }\xi< t_0,


\end{cases}


\end{equation}


%


\begin{equation}


x(t_0)=x_0.


\end{equation}


%


Неоднородность $f$ в уравнении (11) является решением уравнения (6)


при $t_0\ge0$ и некоторой правой части


$p\in\bold L_1[0,+\infty)$, определяемой начальной


функцией $\phi$.





Решение задачи Коши (11), (12) имеет представление (см., напр.,


\cite{Mak})


%


\begin{equation}


x(t)=X(t,t_0)x_0+\int_{t_0}^tX(t,s)f(s)\,ds.


\end{equation}


%


Фундаментальная матрица $X_s(t)=X(t,s)$ является решением матричной


задачи Коши (4) или, как показано в \cite{Mak}, решением матричной задачи


Коши


$$


\Dot X_s(t)=\sum_{g\in G}\sum_{h\in H}R_{gh}X_s(t-g-h),\quad t\ge s\ge t_0,


$$


$$


X_s(t)=0\text{ при }t<s,\qquad X_s(s)=E.


$$





Формула (13) позволяет установить связь равномерной экспоненциальной


устойчивости с экспоненциальной оценкой


%


\begin{equation}


\|X(t,s)\|\le Me^{-\alpha(t-s)},\quad t\ge s\ge t_0\ge0,


\end{equation}


%


где $M,\alpha>0$.





Основным результатом статьи является





\begin{thmnonum} Уравнение {\em (1)} равномерно экспоненциально устойчиво


тогда и только тогда, когда устойчиво уравнение {\em (6)} и фундаментальная


матрица уравнения {\em (11)} имеет экспоненциальную оценку {\em (14)}.


\end{thmnonum}





Для пояснения идеи доказательства теоремы отметим некоторые свойства


уравнения (1).





Как хорошо известно (\cite{Az1}, гл.\,5), решение задачи


%


\begin{equation}


\begin{cases}


\Dot x(t)=\sum\limits_{g\in G_1}A_g(t)\Dot x(t-g)+


\sum\limits_{h\in H}B_h(t)x(t-h)+f(t), & t\ge0 ,\\


%


x(\xi)=0 & \text{при }\xi<0,


\end{cases}


\end{equation}


%


$$


x(0)=x_0


$$


для любого $f\in\bold L_1[0,+\infty)$ имеет представление


%


\begin{equation}


x(t)=X(t,0)x_0+\int_0^tC(t,s)f(s)\,ds.


\end{equation}


%





Формула (16) является распространением на уравнение (15) известной


формулы представления решения задачи Коши для линейного обыкновенного


дифференциального уравнения. Для уравнений, разрешенных относительно


производной, такое представление было изучено в самой общей ситуации в


\cite{Mak}. Роль формулы (16), имеющей место для широкого класса


функционально-дифференциальных уравнений, отмечалась, например, в \cite{Hale},


\cite{Kol}, \cite{Ku3}, \cite{Az1}--\cite{Mal}, \cite{Mak}.


Более глубокое исследование свойств матрицы Коши $C(t,s)$


существенно осложняется неопределенностью в выборе при каждом $t>0$


представителя класса эквивалентности $C(t,\cdot)$. Например, возникают


трудности при изучении вопроса о поведении решений задачи (1), (2),


когда последовательность возмущений $f$ сходится к $\delta$-функции.


Возникают большие трудности с определением матрицы Коши как


обобщенного решения уравнения (1). Поэтому вопрос о наличии


экспоненциальной оценки (5) не совсем корректен. Это неравенство надо


понимать почти всюду. Справедлива





\begin{lemnonum} Уравнение {\em (1)} равномерно экспоненциально устойчиво


тогда


и только тогда, когда найдутся положительные константы $N$ и


$\gamma$ и измеримое множество $\Delta\subset[0,+\infty)$, такие, что


$\text{mes}\,([0,+\infty)\backslash\Delta)=0$


и при всех $t\in[0,+\infty)$, $s\in\Delta$, $t\ge s$


справедлива оценка {\em (5)}.


\end{lemnonum}





\begin{notenonum} При доказательстве теоремы установлена


эквивалентность следующих свойств уравнения (1):


\begin{itemize}


\item[1)] равномерная экспоненциальная устойчивость;


\item[2)] экспоненциальная оценка (5) матрицы Коши $C(t,s)$;


\item[3)] сильная $\bold L_1$-устойчивость (т.е. 


$x,\Dot x\in\bold L_1[0,+\infty)$


при любом


$f\in\bold L_1[0,+\infty)$


в уравнении (15));


\item[4)] обратимость оператора ${\cal I-S}$ в пространстве 


$\bold L_1[0,+\infty)$ и экспоненциальная оценка (14) фундаментальной матрицы


$X(t,s)$.


\end{itemize}


\end{notenonum}





Проведенное в \cite{Az2} изучение сильной устойчивости показывает, что для


ЛДРУ любой признак такой устойчивости (и, следовательно,


экспоненциальной оценки (5) матрицы Коши $C(t,s)$) может быть получен


$W$-методом.





В заключение отметим, что при доказательстве теоремы предлагаемой


работы были использованы результаты Л.\,Вейса (напр., \cite{Weis}) о


представлении порядково непрерывных операторов, действующих в


пространствах измеримых функций. Кроме того, в доказательстве


использована идея, близкая к идее доказательства теоремы Ю.С.\,Колесова


из работы \cite{Kol1}. Разрешимость относительно производной как необходимое


условие экспоненциальной устойчивости по норме $\bold C$


рассматривал В.Г.\,Курбатов в статье \cite{Ku4}. В работе \cite{Ku5}


выделен класс


таких ЛДРУ, для которых $\bold C$-устойчивость, $\bold 


C^1$-устойчивость и $\bold W_\infty^{(1)}$-устойчивость


эквивалентны.





\section{Представление решений ЛДРУ}





Следуя схеме, изложенной в \cite{Az1}, сведем неоднородную задачу


%


\begin{equation}


\Dot x(t)=\sum\limits_{g\in G_1}A_g(t)\Dot x(t-g)+


\sum\limits_{h\in H}B_h(t)x(t-h)+f(t),\quad t\in[0,b],


\end{equation}


%


\begin{equation}


x(\xi)=0\quad\text{при }\xi\le 0


\end{equation}


%


к линейному операторному уравнению в пространстве $\bold L_1[0,b]$.


Обозначая $\Dot x=z$ и учитывая (18), из (17) получим уравнение


%


\begin{equation}


({\cal I-S}_b)z-{\cal Q}_bz=f.


\end{equation}


%





Вольтерровы операторы


${\cal S}_b,{\cal Q}_b:\bold L_1[0,b]\to\bold L_1[0,b]$


действуют по правилам


%


\begin{align*}


({\cal S}_bz)(t) &=\sum_{g\in G_1}A_g(t)({\cal S}_{g,b}z)(t),\\


({\cal Q}_bz)(t) &=\sum_{h\in H}B_h(t)({\cal S}_{h,b}{\cal J}_bz)(t),


\end{align*}


%


где


$$


({\cal S}_{g,b}z)(t)=


\begin{cases}


z(t-g) & \text{ при }g\le t\le b;\\


%


0 &\text{ при }0\le t<g,


\end{cases}


$$


%


$$


({\cal J}_bz)(t)=\int_0^tz(s)\,ds.


$$





Заметим, что нормы $\|{\cal S}_b\|_{1,b}$


и


$\|{\cal Q}_b\|_{1,b}$


операторов ${\cal S}_b,{\cal Q}_b:\bold L_1[0,b]\to\bold L_1[0,b]$


удовлетворяют оценкам


%


\begin{align*}


\|{\cal S}_b\|_{1,b} &\le\sum_{g\in G_1}\vraisup_{t\ge0}\|A_g(t)\|,\\


%


\|{\cal Q}_b\|_{1,b} &\le\sum_{h\in H}\int_0^b\|B_h(t)\|dt.


\end{align*}


%





В условиях основной теоремы оператор $({\cal I-S}_b):\bold 


L_1[0,b]\to\bold L_1[0,b]$ вольтеррово обратим. Обратный оператор


${\cal D}_b=({\cal I-S}_b)^{-1}$ имеет представление


%


\begin{equation}


({\cal D}_by)(t)=y(t)+\sum_{g\in G}P_g(t)({\cal S}_{g,b}y)(t),


\quad t\in[0,b].


\end{equation}


%


Коэффициенты $P_g$ вычисляются по формуле (8). Применив к обеим


частям уравнения (19) оператор ${\cal D}_b$, получим уравнение


%


\begin{equation}


z-{\cal P}_bz={\cal D}_bf,


\end{equation}


%


где ${\cal P}_b={\cal D}_b{\cal Q}_b$.


Согласно (\cite{Az1}, гл.\,1, \S\,1.2, с.\,18),


оператор ${\cal Q}_b$ компактен, что влечет


компактность оператора ${\cal P}_b$. Спектральный радиус компактного


вольтеррова оператора ${\cal P}_b:\bold L_1[0,b]\to\bold L_1[0,b]$


равен нулю (\cite{Ku3}, гл.\,4, \S\,4.4, п.\,4.4.12, с.\,81;


\cite{Zab}, гл.\,V, \S\,6, п.\,6.2,


с.\,153). Отсюда


%


\begin{equation}


({\cal I-P}_b)^{-1}=\sum_{k=0}^\infty{\cal P}_b^k.


\end{equation}


%





Из (20)--(22) находим


$$


({\cal I-S}_b-{\cal Q}_b)^{-1}=({\cal I-P}_b)^{-1}{\cal D}_b=


({\cal I+P}_b({\cal I-P}_b)^{-1}){\cal D}_b=


{\cal D}_b+{\cal P}_b({\cal I-P}_b)^{-1}{\cal D}_b.


$$





Таким образом, решение $z$ уравнения (19) имеет представление


%


\begin{equation}


z={\cal U}_bf,


\end{equation}


%


где


%


\begin{equation}


{\cal U}_b={\cal D}_b+{\cal K}_b.


\end{equation}


%


Оператор ${\cal K}_b={\cal P}_b({\cal I-P}_b)^{-1}{\cal D}_b$


компактен в


силу компактности оператора ${\cal P}_b$. Тогда, очевидно, оператор


${\cal K}_b$ слабо компактен и согласно замечанию из (\cite{Dun},


гл.\,VI, 8.11, с.\,547)


оператор ${\cal K}_b$ является интегральным оператором: существует


матрица-функция $K_b:[0,b]\times[0,b]\to\bold M^n$ такая, что


%


\begin{equation}


({\cal K}_by)(t)=\int_0^bK_b(t,s)y(s)\,ds.


\end{equation}


%


Так как оператор ${\cal K}_b$ вольтерров, то $K_b(t,s)=0$ для каждого


$t\in[0,b]$ при почти всех $s>t$.





Решение $x=\col\{x_1,\dotsc,x_n\}$ задачи (17), (18) найдем, интегрируя


равенство (23)


\begin{gather*}


x_i(t)=\biggl\langle\int_0^tz(s)ds,e_i\biggr\rangle =


\biggl\langle\int_0^t({\cal U}_bf)(s)ds,e_i\biggr\rangle= 


\int_0^b\langle f(s),({\cal U}_b^*\chi_{[0,t]}e_i)(s)\rangle ds=\\


%


=\int_0^b\biggl\langle\sum_{j=1}^n \langle f(s),


e_j\rangle e_j,({\cal U}_b^*\chi_{[0,t]}e_i)(s)\biggr\rangle ds= 


\int_0^b\sum_{j=1}^nf_j(s)\langle e_j,({\cal U}_b^*\chi_{[0,


t]}e_i)(s)\rangle ds.


\end{gather*}





Таким образом, решение задачи (17), (18) допускает интегральное


представление


%


\begin{equation}


x(t)=\int_0^tC_b(t,s)f(s)\,ds


\end{equation}


%


с ядром $C_b(\cdot,\cdot)$, определяемым равенством


%


\begin{equation}


C_{bij}(t,s)=(\kappa\langle e_j,({\cal U}_b^*\chi_{[0,t]}e_i)(\cdot)\rangle )(s).


\end{equation}


%


Здесь $\kappa$ -- любой оператор выбора представителя из класса


эквивалентных по мере Лебега измеримых и ограниченных в существенном


функций, применяемый к скалярной функции


$\langle e_j,({\cal U}_b^*\chi_{[0, t]}e_i)(\cdot)\rangle$.





Продолжим изучение свойств интегрального представления (26). Из


(24) следует равенство


%


\begin{equation}


{\cal U}_b^*={\cal D}_b^*+{\cal K}_b^*.


\end{equation}


%


Опуская простые вычисления, основанные на представлениях (20) и (25),


приведем вид сопряженных операторов


%


\begin{equation}


({\cal D}_b^*v)(s)=v(s)+\sum_{g\in G}R_g^\top (s+g)({\cal S}_{-g,


b}v)(s),


\end{equation}


%


\begin{equation}


({\cal K}_b^*v)(s)=\int_s^bK_b^\top (t,s)v(t)dt,


\end{equation}


%


где


$$


({\cal S}_{g,b}^*v)(s)=({\cal S}_{-g,b}v)(s)=


\begin{cases}


z(t+g) & \text{ при }0\le t\le b-g;\\


%


0 & \text{ при }b-g<t\le b,


\end{cases}


$$


%


{``}${}^\top ${''} --- знак транспонирования матрицы.





Полученные равенства запишем в виде интегралов по $\bold M^n$-значной


(матричной) случайной мере \cite{Weis}. Определим случайные меры 


$\mu_b^a(s,\Delta)$


и


$\mu_b^i(s,\Delta)$ борелевского множества


$\Delta\subset[0,b]$ равенствами


%


\begin{equation}


\mu_b^a(s,\Delta)=\sum_{g\in G}\chi_\Delta(s+g)P_g^\top(s+g)+


\chi_\Delta(s)


\end{equation}


%


и


%


\begin{equation}


\mu_b^i(s,\Delta)=\int_\Delta K_b^\top (\vartheta,s)\,d\vartheta.


\end{equation}


%





Отметим, что выше неявно и неконструктивно произведен выбор


представителей из классов эквивалентности коэффициентов $P_g^\top$ и


$K_b^\top (t,\cdot)$.





Меры, определенные равенствами (31) и (32), обладают всеми свойствами


случайных мер: для каждого борелевского множества $\Delta\subset[0,b]$


функции $\mu_b^a(\cdot,\Delta)$


и


$\mu_b^i(\cdot,\Delta)$


измеримы, при каждом $s\in[0,b]$ функции множества $\mu_b^a(s,\cdot)$


и $\mu_b^i(s,\cdot)$ -- борелевские меры. Для любой борелевской


вектор-функции $v:[0,b]\to\bold R^n$ равенство (29) есть случайный


интеграл


$$


({\cal D}_b^*v)(s)=\int_{[0,b]}\mu_b^a(s,d\vartheta)v(\vartheta),


$$


а равенство (30) --- случайный интеграл


$$


({\cal K}_b^*v)(s)=\int_{[0,b]}\mu_b^i(s,d\vartheta)v(\vartheta).


$$





Классы эквивалентности этих случайных интегралов не зависят от


конкретного выбора борелевских представителей класса 


$v\in\bold L_\infty[0,b]$.





Согласно (28) имеем


%


\begin{equation}


({\cal U}_b^*v)(s)=\int_{[0,b]}\mu_b(s,d\vartheta)v(\vartheta),


\end{equation}


%


где $\mu_b=\mu_b^a+\mu_b^i$. В силу (31) компонента $\mu_b^a$ является


дискретной (атомарной) случайной мерой, а в силу (32) компонента


$\mu_a^i$ является случайной мерой, абсолютно непрерывной относительно


меры Лебега $\mathrm{mes}$. Формула (33) позволяет ``конструктивно'' определить


при каждом $t\in[0,b]$ процедуру выбора представителя из класса


эквивалентности ядра $C_b(t,\cdot)$ в интегральном представлении (26).


Равенство (27) определяет случайный интеграл


%


\begin{equation}


C_b(t,s)=\int_{[0,b]}\chi_{[0,t]}(\vartheta)\mu_b(s,d\vartheta)


\end{equation}


%


или, что эквивалентно, --- случайную меру


$$


C_b(t,s)=\mu_b(s,[0,t]).


$$





Таким образом, при каждом $s\in[0,b]$ функция $C_b(\cdot,s)$ является


функцией распределения матричной меры $\mu_b(s,\cdot)$. Эта функция


имеет ограниченную на $[0,b]$ вариацию и непрерывна справа. Кроме


того, $C_b(\cdot,s)=C_b^a(\cdot,s)+C_b^i(\cdot,s)$, так как


$\mu_b=\mu_b^a+\mu_b^i$. Компонента $C_b^a(\cdot,s)$ является функцией


распределения дискретной меры. Поэтому при каждом $s\in[0,b]\;$


$\;C_b^a(\cdot,s)$ --- ступенчатая функция, имеющая разрывы в точках


$t=s+g$, $g\in G$.





Согласно (31)


%


\begin{equation}


C_b(s+g+0,s)-C_b(s+g-0,s)=P_g^\top (s,g).


\end{equation}


%





Компонента $C_b^i(\cdot,s)$ при каждом $s\in[0,b]$ является функцией


распределения абсолютно непрерывной меры. Поэтому функция


$C_b^i(\cdot,s)$ абсолютно непрерывна. Свойства разложения 


$C_b(\cdot,s)=C_b^a(\cdot,s)+C_b^i(\cdot,s)$


существенно используются при


доказательстве теоремы. Особо подчеркиваем, что ступенчатая


компонента $C_b^a(\cdot,s)$ определяется разностным оператором


${\cal D}_b^*$.





\section{Экспоненциальная оценка матрицы Коши. Доказательство леммы}





Как отмечено в п.\,1, решение задачи Коши (15) для любого


$f\in\bold L_1[0,+\infty)$ имеет представление в виде формулы (16).





Матрица Коши $C(t,s)$ определяется своими сужениями $C_b(t,s)$:


найдется матрица-функция


$C:[0,+\infty)\times[0,+\infty)\to\bold M^n$


такая,


что при любом $b>0$ имеем $C(t,s)=C_b(t,s)$ для каждого $t\in[0,b]$


при почти всех $s\in[0,b]$.





Для определения ядра $C(t,s)$ представления (16) положим


$C(t,s)=C_k(t,s)$ при всех $t\in[k,k+1)$ и почти всех $s\in[0,k+1)$,


где $C_k(t,s)$ для каждого $k=0,1,\dots$ определяется формулой (34).


Тогда при любом $s\in[0,+\infty)$ функция $C(\cdot,s)$ является


непрерывной справа


функцией ограниченной вариации на каждом отрезке $[0,b]$.





{\bf Доказательство леммы. Необходимость.} Пусть $s\in[0,+\infty)$. Для


каждого $y\in\bold R^n$ определим $\delta$-образную последовательность


возмущений $f_k$ равенством


$$


f_k(t)=k\chi_{\Delta_k}(t)y,\quad k=1,2,\dotsc,


$$


где $\Delta_k=[s,s+1/k]$.





Решение $x$ задачи (15) при $f=f_k$ есть решение задачи (1), (2) при


$t_0=s+1/k$ и начальной функции


$\phi_k\in\bold W_1(-\infty,t_0]$,


являющейся сужением функции $x_k$ на полуось $(-\infty,t_0]$. В


силу (16) имеем $\phi(t)=0$ при $t\le s$.





Из ограничений на коэффициенты уравнения (1) следует существование


такого положительного $c$ (см. \cite{Ku3}, гл.\,4, \S\,4.1, п.\,4.1.5,


с.\,73), что


$$


\|\phi_k\|_{\bold W_1(-\infty,t_0]}\le c\|f_k\|_{\bold L_1[0,+\infty)}.


$$


Из условий теоремы имеем


$$


\|x_k(t)\|\le


N_1e^{-\gamma(t-t_0)}\|\phi_k\|_{\bold W_1(-\infty,t_0]}.


$$


Отсюда с учетом $x_k(t)=0$ при $t\le s$ получаем


%


\begin{equation}


\|x_k(t)\|\le Ne^{-\gamma(t-s)}\|f_k\|_{\bold L_1[0,+\infty)}\le


Ne^{-\gamma(t-s)}\|y\|


\end{equation}


%


при $s\le t\le s+1/k$, где $N=c\exp\gamma\cdot\max\{1,N_1\}$.





Из формулы Коши (16) получим


$$


x_k(t)=\int_s^tC(t,\vartheta)f_k(\vartheta)\,d\vartheta=


k\int_s^{s+1/k}C(t,\vartheta)y\,d\vartheta.


$$


Согласно теореме Лебега о предельном переходе под знаком интеграла


(\cite{Dun}, гл.\,III, 6.16, с.\,168) при каждом $t\in[0,+\infty)$


последовательность $x_k(t)$ сходится при почти всех $s\in[0,+\infty)$ к


$C(t,s)y$. Отсюда и из (36) заключаем следующее: при фиксированных


$t\in[0,+\infty)$ и $y\in\bold R^n$ найдется измеримое множество


$\Delta(t,y)\subset[0,+\infty)$


такое, что


$\mathrm{mes}\,([0,+\infty)\backslash\Delta(t,y))=0$


и при всех $s\in\Delta(t,y)$ справедлива оценка


%


\begin{equation}


\|C(t,s)y\|\le Ne^{-\gamma(t-s)}\|y\|.


\end{equation}


%





Выберем счетное всюду плотное в $[0,+\infty)$ подмножество ${\cal M}$ и


счетное всюду плотное в $\bold R^n$ подмножество ${\cal Y}$. Множество


$\Delta=\bigcap\limits_{t\in{\cal M}}\bigcap\limits_{y\in{\cal


Y}}\Delta(t,y)$


является множеством полной меры как пересечение


счетного набора множеств полной меры. При всех $t\in{\cal M}$ и


$s\in\Delta$ неравенство (37) выполнено на плотном подмножестве 


${\cal Y}\subset\bold R^n$. Поэтому из (37) следует при всех $t\in{\cal M}$ и


$s\in\Delta$ оценка (5).





Для фиксированного $t\in[0,+\infty)$ подберем монотонно убывающую и


сходящуюся к $t_0$ последовательность $\{t_k\}\subset{\cal M}$. В силу


непрерывности справа матрицы-функции $C(\cdot,s)$ получим


$$


\|C(t,s)\|=\lim_{k\to\infty}\|C(t_k,s)\|\le Ne^{-\gamma(t-s)}.


$$





{\bf Достаточность.} Пусть $\delta=\max\{g,h\}(g\in G_1,h\in H)$.


Решение $x$ задачи (1), (2) с начальными данными $t_0=s\ge0$,


$\phi\in\bold W_1(-\infty,s]$ имеет при $s\le t\le s+\delta$


оценку


$$


\|x(t)\|\le c_1\delta\|\phi\|_{\bold W_1(-\infty,s]}.


$$





Положим $\rho=\min\{1,s\}$. Определим функцию $\psi$ условиями:


$\psi(t)=x(t)$ при $t\in[s,s+\rho]$,


$\psi(t)=x(s)(1+\dfrac{t-s}{\rho})$


при $t\in[s-\rho,s]$


и


$\psi(t)=0$ 


при


$t\le s-\rho$. Решение задачи (1), (2) при всех $t\ge s$ совпадает с


решением задачи (15) при возмущении $f$, определяемом равенством


$$


f(t)=


\begin{cases}


0 & \text{ при } t\ge s+\delta;\\


%


\Dot\psi(t)-\sum\limits_{g\in G_1}A_g\Dot\psi(t-g)-\sum\limits_{h\in


H}B_h(t)\psi(t-h) & \text{ при }t<s+\delta.


\end{cases}


$$


%


Согласно построению


$$


\|f\|_{\bold L_1[0,+\infty)}\le c\|x\|_{\bold W_1[s,s+\delta]}\le


c_1c\delta\|\phi\|_{\bold W_1(-\infty,s]}.


$$


В предположении оценки (5) при $t\ge s+\delta$ имеем


%


\begin{multline*}


\|x(t)\|\le N\int_0^{s+


\delta}e^{-\gamma(t-\vartheta)}\|f(\vartheta)\|d\vartheta\le


Ne^{\gamma s}e^{-\gamma(t-s)}\int_0^{s+\delta}e^{


-\gamma(s+\delta-\vartheta)}\|f(\vartheta)\|d\vartheta\le \\


%


\le Ne^{\gamma


s}e^{-\gamma(t-s)}\int_0^{+\infty}\|f(\vartheta)\|d\vartheta


=Ne^{\gamma s}e^{-\gamma(t-s)}\|f\|_{\bold L_1[0,+\infty)}\le 


Nc_1c\delta e^{\gamma


s}e^{-\gamma(t-s)}\|\phi\|_{\bold W_1(-\infty,s]}.


\end{multline*}


Лемма доказана.





\section{Доказательство основной теоремы}





{\bf Необходимость.} Пусть уравнение (1) равномерно экспоненциально


устойчиво. Тогда в силу леммы матрица Коши $C(t,s)$ уравнения (15)


имеет экспоненциальную оценку (5). Из результатов п.\,2 следует,


что при фиксированном $s\in\Delta$ точки разрыва функции $C(\cdot,s)$


располагаются на дискретном множестве $G(s)=\{s+h:g\in G\}$.





Из (35) и (5) при каждом $s\in\Delta$ имеем


%


\begin{multline}


\|P_g^\top(s+g)\|=\|C(s+g+0,s)-C(s+g-0,s)\|\le \\


\le \|C(s+g+0,s)\|+\|C(s+g-0,s)\|\le2Ne^{-g}.


\end{multline}


%





Число слов длины $k$ в коммутативной полугруппе $G\subset[0,+\infty)$,


порожденной конечным множеством $G_1=\{g_1,\dotsc,g_l\}$, не


превосходит $\sigma k^{l-1}$, $\sigma>0$. Число $g$, соответствующее


слову длины $k=0,1,\dotsc$, не меньше, чем $kd$, где 


$d=\min\{g:g\in G_1\}>0$, поэтому из (38) следует


$$


\|P_g^\top(s+g)\|\le2Ne^{-kd},


$$


$$


\sum_{g\in G}\|P_g^\top(s+g)\|=\sum_{k=1}^\infty


\sum\begin{Sb}


g=g_1+\dotsb+g_k,\\


g_k\in G_1


\end{Sb}


\|P_g^\top(s+g)\|\le2N\sigma\sum_{k=1}^\infty k^{l-1}e^{-kd}<\infty.


$$





Из равенства (29) для борелевской вектор-функции $v\in\bold 


L_\infty[0,b]$ следует, что для всех $s\in\Delta$ имеют место


неравенства


\begin{multline*}


\|({\cal D}_b^*v)(s)\|\le\|v(s)\|+\sum_{g\in G}\|P_g^\top(s+


g)\|\cdot\|({\cal S}_{-g,b}v)(s)\|\le \\


%


\le\|v(s)\|+\biggl(\sup\limits_{s\in\Delta}\biggl\{\sum_{g\in


G}\|P_g^\top(s+g)\|\biggr\}\biggr)\biggl(\sup\limits_{g\in


G}\{\|({\cal S}_{-g,b}v)(s)\|\}\biggr)\le \\


%


\le q\cdot\max\Bigl\{\|v(s)\|,\sup\limits_{g\in


G}\{\|({\cal S}_{-g,b}v)(s)\|\}\Bigr\},


\end{multline*}


%


где


$$


q=1+\sup_{s\in\Delta}\sum_{g\in G}\|P_g^\top(s+g)\|.


$$





Так как $\Delta$ -- множество полной меры, то отсюда следует оценка


$$


\|{\cal D}_b^*v\|_{\infty,b}\le q\max\{\|v\|_{\infty,b},


\sup\limits_{g\in G}\{\|{\cal S}_{-g,b}v\|_{\infty,b}\}\}\le q\|v\|_{\infty,b}.


$$


Тогда


$$


\|{\cal D}_b\|_{1,b}=\|{\cal D}_b^*\|_{\infty,b}\le q.


$$


Здесь и ниже $\|v\|_{\infty,b}$ ($\|{\cal Q}\|_{\infty,b}$) --- норма элемента


$v\in\bold L_\infty[0,b]$


(оператора ${\cal Q}:\bold L_\infty[0,b]\to\bold L_\infty[0,b]$),


$\|u\|_{1,b}$ ($\|{\cal P}\|_{1,b}$) ---


норма элемента $u\in\bold L_1[0,b]$ (оператора 


${\cal P}:\bold L_1[0,b]\to\bold L_1[0,b]$).





Итак, получена равномерная по $b>0$ оценка нормы разностного оператора


${\cal D}_b=({\cal I-S}_b)^{-1}:\bold L_1[0,b]\to\bold L_1[0,b]$.


Такая оценка эквивалентна оценке (9) для всех $b>0$. А это и означает


устойчивость разностного уравнения (6). Необходимость первого условия


теоремы доказана.





Задача (15) эквивалентна операторному уравнению


%


\begin{equation}


({\cal I-S})\Dot x-{\cal R}x=f,


\end{equation}


%


где операторы ${\cal R,S}:\bold L_1[0,+\infty)\to\bold L_1[0,


+\infty)$ определены равенствами


%


\begin{align*}


({\cal R}x)(t) &=\sum_{h\in H}B_h(t)({\cal S}_hx)(t),\\


({\cal S}z)(t) &=\sum_{g\in G_1}A_g(t)({\cal S}_gz)(t).


\end{align*}





Из доказательства необходимости первого условия теоремы следует


обратимость оператора ${\cal I-S}$, причем для оператора


${\cal D}=({\cal I-S})^{-1}$ имеем $\|{\cal D}\|_{1,\infty}\le q.$





Так как матрица-функция Коши уравнения (1) согласно лемме имеет


экспоненциальную оценку (5), то справедливы неравенства


%


\begin{multline*}


\|x\|_{1,\infty}=\int_0^{+\infty}\biggl\|\int_0^tC(t,


\vartheta)f(\vartheta)d\vartheta\biggr\|dt\le


N\int_0^{+\infty}\int_0^te^{-\gamma(t-\vartheta)}\|


f(\vartheta)\|d\vartheta\,dt\le \\


%


\le N\int_0^{+\infty}\int_\vartheta^{+\infty}e^{-\gamma(t-\vartheta)}


dt\|f(\vartheta)\|d\vartheta\le


N\int_0^{+\infty}\int_0^{+\infty}e^{-\gamma(t-\vartheta)}


dt\|f(\vartheta)\|d\vartheta\le\frac{N}{\gamma}\|f\|_{1,\infty}.


\end{multline*}





Итак, при любой $f\in\bold L_1[0,+\infty)$ решение $x$ уравнения


(39) суммируемо на $[0,+\infty)$.





Так как $\Dot x={\cal D}({\cal R}x+f)$,


то


$\Dot x\in\bold L_1[0,+\infty)$. Таким образом, при любой


$f\in\bold L_1[0,+\infty)$


имеем $x,\Dot x\in\bold L_1[0,+\infty)$. В терминологии статьи


\cite{Az2}


это свойство означает сильную $\bold L_1$-устойчивость уравнения


(39). Операторы ${\cal R}$ и ${\cal S}$ ввиду условий на коэффициенты


и ограниченности запаздываний удовлетворяют $\Delta\bold 


L_1$-условию \cite{Az2} (или условию экспоненциального убывания памяти


(\cite{Ku1}--\cite{Ku3}, \cite{Si3}, \cite{Si4}, \cite{Ku4}, \cite{Ku5}).


Фундаментальная матрица $X_s(t)=X(t,s)$ уравнения (11) является


(\cite{Az1}, гл.\,3, \S\,3.4, с.\,62; \cite{Mak})


решением матричной задачи Коши


(4) и совпадает с фундаментальной матрицей уравнения (1).





Тогда согласно теореме 4 статьи \cite{Az2} сильная 


$\bold L_1$-устойчивость при выполнении $\Delta\bold L_1$-условия


влечет при некоторых положительных $M$ и $\alpha$ экспоненциальную


оценку (14) для матрицы $X(t,s)$.





{\bf Достаточность.} Будем предполагать, что оператор 


$({\cal I-S}):\bold L_1[0,+\infty)\to\bold L_1[0,+\infty)$


вольтеррово обратим и


фундаментальная матрица $X(t,s)$ имеет экспоненциальную оценку (14).





В силу вольтерровой обратимости оператора 


$({\cal I-S}):\bold L_1[0,+\infty)\to\bold L_1[0,+\infty)$


решение $x$ уравнения (39) при


любой $f\in\bold L_1[0,+\infty)$ имеет представление \cite{Mak}


$$


x(t)=\int_0^tX(t,s)({\cal D}f)(s)ds.


$$





Используя оценку (14) для $X(t,s)$ и рассуждая так же, как в п.\,1,


получим неравенства


$$


\|x\|_{1,\infty}\le\frac{M}{\alpha}\|{\cal D}\|_{1,\infty}\cdot\|f\|_{1,


\infty}\le\frac{M}{\alpha}q\|f\|_{1,\infty}.


$$





Итак, $x\in\bold L_1[0,+\infty)$


при любой


$f\in\bold L_1[0,+\infty)$.


Далее из равенства $\Dot x={\cal D}({\cal R}x+f)$ получаем, что


$\Dot x\in\bold L_1[0,+\infty)$. Таким образом, уравнение (39)


сильно $\bold L_1$-устойчиво. Отсюда согласно теореме 5 из \cite{Az2}


следует экспоненциальная оценка (5) матрицы Коши $C(t,s)$ и в силу


леммы --- равномерная экспоненциальная устойчивость. $\square$








\begin{thebibliography}{99} 


 


\bibitem{Hale} 


Хейл Дж. {\em Теория функционально-дифференциальных уравнений.}


-- М.: Мир, 1984. -- 424 с.





\bibitem{Kol} 


Колмановский В.Б., Носов В.Р. {\em Устойчивость и периодические


режимы регулируемых систем с последействием.} -- М.: Наука, 1981. --


448~с.





\bibitem{Ku1} 


Курбатов В.Г. {\em О разрешимости относительно производной


устойчивого функционально-дифференциального уравнения} // Укр. матем.


журн. -- 1982. -- Т.\,34. -- \No\,1. -- С.\,103--106.





\bibitem{Ku2} 


Курбатов В.Г. {\em О сведении устойчивого уравнения нейтрального типа


к уравнению запаздывающего типа} // Сиб. матем. журн. -- 1983. -- Т.\,24.


-- \No\,4. -- С.\,209--212.





\bibitem{Ku3} 


Курбатов В.Г. {\em Линейные дифференциально-разностные


уравнения.} -- Воронеж: Изд-во ВГУ, 1990. -- 168~с.





\bibitem{Si1} 


Симонов П.М., Чистяков А.В. {\em О необходимом условии


устойчивости дифференциально-разностных уравнений нейтрального типа}


// Моделирование и исследование устойчивости систем. Тез. докл. конф.


Ч.2. -- Киев, 1993. -- С.\,39.





\bibitem{Si2} 


Симонов П.М., Чистяков А.В. {\em Об экспоненциальной


устойчивости линейных дифференциально-разностных уравнений,


неразрешенных относительно производной} // Вестн. ПГТУ. Сер. матем. и


прикл. матем. -- Пермь, 1996. -- \No\,1. -- С.\,44--49.





\bibitem{Az1} 


Азбелев Н.В., Максимов В.П., Рахматуллина Л.Ф. {\em Введение в теорию


функционально-дифференциальных уравнений.} -- М.: Наука, 1991. --


280~с.





\bibitem{Az2} 


Азбелев Н.В., Березанский Л.М., Симонов П.М., Чистяков А.В.


{\em Устойчивость линейных систем с последействием.} IV // Дифференц.


уравнения. -- 1993. -- Т.\,29. -- \No\,2. -- С.\,196--204.





\bibitem{Ber} 


Березанский Л.М., Малыгина В.В., Соколов В.А. {\em Признаки


экспоненциальной устойчивости решений уравнений с ограниченным


последействием} // ДАН СССР. -- 1986. -- Т.\,289. -- \No\,1. -- С.\,11--14.





\bibitem{Mal} 


Малыгина В.В. {\em Из истории развития теории устойчивости уравнений


с постоянным запаздыванием} // Функц.-дифференц. уравнения. -- Пермь,


1991. -- С.\,70--78.





\bibitem{Si3} 


Симонов П.М., Чистяков А.В. {\em Об обратимости линейных операторов,


коммутирующих со сдвигом} // Функц.-дифференц. уравнения. -- Пермь, 1992.


-- С.\,144--159.





\bibitem{Si4} 


Симонов П.М., Чистяков А.В. {\em О разрешимости периодических


уравнений} // Вестн. ПГТУ. Сер. матем. и прикл. матем. --


Пермь, 1994. -- \No\,1. -- С.\,61--71.





\bibitem{Mak} 


Максимов В.П. {\em О формуле Коши для функционально-дифференциального


уравнения} // Дифференц. уравнения. -- 1977. -- Т.\,13. --


\No\,4. -- С.\,601--606.





\bibitem{Weis} 


Weis L. {\em Decompositions of positive operators and some of their


applications} // Funct. Anal.: Surv. and Recent Results. 3: Proc. 3rd


Conf. -- Paderborn. -- Amsterdam e.a., 1984. -- P.\,95--115.





\bibitem{Kol1} 


Колесов Ю.С. {\em Об устойчивости решений линейных


дифференциально-разностных уравнений нейтрального типа} // Сиб. матем.


журн. -- 1979. -- Т.\,20. -- \No\,2. -- С.\,317--321.





\bibitem{Ku4}


Курбатов В.Г. {\em Об уравнениях нейтрального типа с внешним


дифференцированием} // Дифференц. уравнения. -- 1987. -- Т.\,23. -- \No\,3.


-- С.\,434--442.





\bibitem{Ku5} 


Курбатов В.Г. {\em Об устойчивости дифференциально-разностных


уравнений с производной и без производной} // Дифференц. уравнения.


-- 1988. -- Т.\,24. -- \No\,9. -- С.\,1503--1509.





\bibitem{Zab} 


Забрейко П.П., Кошелев А.И., Красносельский М.А. и


др. {\em Интегральные уравнения.} -- М.: Наука, 1968. -- 448~с.





\bibitem{Dun} 


Данфорд Н., Шварц Дж.Т. {\em Линейные операторы. Общая теория.}


-- М.: Ин.~лит., 1962. -- 895~с.





\end{thebibliography} 


 


\vskip 2cm 





\post{\it Пермский государственный}{\it Поступила}


\post{\it университет}{28.06.1996} 


 


\end{document} 


